\subsection{不可约多项式}

定义2. --- 设 K 是一个交换域。我们说 \(f \in \mathbf{K}[\mathbf{X}]\) 是不可约的，如果 \(\deg f \geqslant 1\) 且 \(f\) 不能被任何多项式 \(g\) （满足 \(0 < \deg g < \deg f\) ）整除。

这等同于说 \(\deg f \geqslant 1\) 且 f 在 K{[}X{]} 中的因子只有标量 \(\neq 0\) 以及 f 与标量 \(\neq 0\) 的乘积。由于关系式 \((f) \subset (g)\) 意味着 g 整除 f，我们看到 K{[}X{]} 的不可约多项式也可以定义为使理想 \((f)\) 为极大理想的多项式 f（I，第 104 页）。

设 \(f, g \in \mathbf{K}[\mathbf{X}]\) 。若 \(f\) 不可约，显然要么 \(f\) 与 \(g\) 互素，要么 \(f\) 整除 \(g\) 。若 \(f\) 与 \(g\) 都不可约，则要么 \(f\) 与 \(g\) 互素，要么每一个都是另一个与标量 \(\neq 0\) 的乘积。特别地，两个不同的不可约首一多项式是互素的。

命题13. —— 设 \(\mathcal{I}\) 为 K{[}X{]} 中不可约首一多项式的集合。设 f 为 K{[}X{]} 的非零元素，且 α 为其首项系数；则存在唯一一族具有有限支撑的正整数 \((v_{p})_{p\in \mathcal{I}}\) ，使得我们有分解

\[
f = \alpha \prod_ {p \in \mathcal {I}} p ^ {v _ {p}}.\tag{18}
\]

只需在 \(f\) 是首一的（即 \(\alpha = 1\) ）情形下证明该命题。我们将对次数 \(n\) （ \(f\) 的）作归纳， \(n = 0\) 的情形是平凡的。于是设 \(n \geqslant 1\) ，并设该命题已对所有次数 \(< n\) 的多项式建立。

设 E 为整除 f 的首一多项式 \(\neq1\) 组成的集合；我们有 \(f\in E\) ，因此 E 非空，且 E 中存在一个次数最小的多项式 g。显然 g 是不可约的，并且存在一个次数 \textless n 的首一多项式 h 使得 f=gh；根据归纳假设，h 是有限族不可约首一多项式的乘积，因此 f 也具有同样的性质。这证明了分解式 (18) 的存在性。

现在让我们证明分解式 (18) 的唯一性。设 \((w_{p})_{p\in \mathcal{I}}\) 为一个具有有限支撑的正整数族，使得 \(f = \prod_{p\in \mathcal{I}}p^{w_p}\) 。由于 \(f\) 的

次数 \(n \geqslant 1\) ，存在 \(p \in \mathcal{I}\) 使得 \(w_{p} > 0\) ；如果我们有 \(v_{p} = 0\) ，则 f 将是 \(\mathcal{I}\) 中异于 p 的元素组成的族之乘积，因此它与 p 互素（IV，第 13 页，推论 6），这与 p 整除 f 的事实相矛盾。由归纳假设，多项式 f/p 有唯一的类型 (18) 的分解；因此我们得出等式 \(w_{q} = v_{q}\) （对每个 \(q \in \mathcal{I}\) ）。

设 \(f\) 是 \(\mathbf{K}[\mathbf{X}]\) 中的一个非零多项式。我们就说 \(f\) 无重因子，如果分解 (18) 中的指数 \(v_{p}\) 全都 \(\leqslant 1\) ；这等同于说 \(f\) 是有限个两两不同的不可约多项式之乘积，或者说 \(f\) 不能被 \(\mathbf{K}[\mathbf{X}]\) 中任何非常数多项式的平方整除。

